% ======================================================================
% Week 4: FLUID DYNAMICS APPLICATIONS
% ======================================================================
\section{Week 4 - Applications of Bernoulli \& Conservation Laws}

\subsection{Rectilinear (or Nearly) Streamlines}

Consider a fluid element moving along a rectilinear streamline. If we neglect dynamic pressure variations transverse to the flow, we recover the hydrostatic pressure distribution:
\begin{equation*}
    p_{1} + \rho g z_{1} = p_{2} + \rho g z_{2}
\end{equation*}
\begin{equation}
    p_{2} - p_{1} = -\rho g (z_{2} - z_{1}) = -\rho g h \implies p_{1} = p_{2} + \rho g h
\end{equation}

\subsection{Example: Pitot-Tube (Speed Sensor)}

A Pitot tube is used to measure the local velocity of the fluid by comparing static and stagnation pressures.

\begin{center}
\begin{tikzpicture}[scale=1.2]
    % Flow field
    \draw[->, thick, primary] (-2, 1) -- (-0.5, 1);
    \draw[->, thick, primary] (-2, 0.5) -- (-0.5, 0.5);
    \draw[->, thick, primary] (-2, 0) -- (-0.5, 0) node[midway, above] {$v$};
    
    % Pitot tube structure
    \draw[thick, fill=gray!20] (0, 0.2) -- (3, 0.2) -- (3, 2) -- (3.4, 2) -- (3.4, -0.2) -- (0, -0.2) arc (-90:90:0.2);
    \filldraw[black] (0,0) circle (1.5pt) node[left] {$2$};
    \filldraw[black] (-1.5,0) circle (1.5pt) node[below] {$1$};
    
    % Streamline
    \draw[thick, dashed, accent] (-3, 0) -- (0, 0) node[midway, below] {Streamline};
\end{tikzpicture}
\end{center}

\textbf{Task:} Determine the free-stream velocity $v$.
Along the streamline connecting point 1 (free stream) and point 2 (stagnation point), we apply Bernoulli's equation, neglecting gravity ($z_1 \approx z_2$):
\[
p_{1} + \frac{1}{2}\rho v_{1}^{2} = p_{2} + \frac{1}{2}\rho v_{2}^{2}
\]
Since point 2 is a stagnation point, $v_2 = 0$. The free stream velocity is $v_1 = v$.
\[
p_{1} + \frac{1}{2}\rho v^{2} = p_{2} + 0 \implies p_{2} - p_{1} = \frac{1}{2}\rho v^{2}
\]

\begin{resultboxnamed}{Velocity from Pitot Tube}
\[
v = \sqrt{\frac{2(p_{2} - p_{1})}{\rho}}
\]
\end{resultboxnamed}

\subsection{Mass Conservation}

Consider a stream tube with varying cross-sectional area.
\begin{center}
\begin{tikzpicture}
    % Stream tube
    \draw[thick] (0, 1) .. controls (2, 0.5) and (4, 1.5) .. (6, 2);
    \draw[thick] (0, -1) .. controls (2, -0.5) and (4, -1.5) .. (6, -2);
    
    % Area 1
    \draw[thick, fill=primary!20] (0,0) ellipse (0.3 and 1);
    \node at (0, 1.3) {$A_1$};
    \draw[->, thick, accent] (-1, 0) -- (0, 0) node[above left] {$v_1, \rho_1$};
    
    % Area 2
    \draw[thick, fill=primary!20] (6,0) ellipse (0.5 and 2);
    \node at (6, 2.3) {$A_2$};
    \draw[->, thick, accent] (6, 0) -- (7.5, 0) node[above right] {$v_2, \rho_2$};
\end{tikzpicture}
\end{center}

\begin{itemize}
    \item \textbf{Mass flow rate} ($\dot{m}$): $[\text{kg/s}]$. $\dot{m} = \rho v A$
    \item \textbf{Volume flow rate} ($\dot{V}$): $[\text{m}^3/\text{s}]$. $\dot{V} = v A$
\end{itemize}

For steady flow through a stream tube, mass is conserved:
\[
\rho_{1}v_{1}A_{1} = \rho_{2}v_{2}A_{2}
\]
Assuming steady and \textbf{incompressible} flow ($\rho_1 = \rho_2$):
\begin{eqbox}{Volume Conservation}
\[
v_{1}A_{1} = v_{2}A_{2}
\]
\end{eqbox}

\subsection{Example: Venturi Meter}

A Venturi meter measures the flow rate in a pipe by constricting the flow and measuring the pressure drop.

\begin{center}
\begin{tikzpicture}[scale=1]
    % Pipe shape
    \draw[thick] (0, 1.5) -- (2, 1.5) -- (3, 0.6) -- (5, 0.6) -- (6, 1.5) -- (8, 1.5);
    \draw[thick] (0, -1.5) -- (2, -1.5) -- (3, -0.6) -- (5, -0.6) -- (6, -1.5) -- (8, -1.5);
    
    % Diameters
    \draw[<->] (1, -1.4) -- (1, 1.4) node[midway, fill=white] {$D$};
    \draw[<->] (4, -0.5) -- (4, 0.5) node[midway, right] {$d$};
    
    % Points
    \filldraw[black] (1, 0) circle (1.5pt) node[below] {$1$};
    \filldraw[black] (4, 0) circle (1.5pt) node[below] {$2$};
    
    % Streamline
    \draw[dashed, accent, thick, ->] (0,0) -- (1,0) .. controls (2,0) and (3,0) .. (4,0) -- (8,0);
\end{tikzpicture}
\end{center}

\textbf{Question:} What is the flow rate?
Applying Bernoulli's equation along the central streamline (assuming horizontal flow, $z_1 = z_2$):
\[
\frac{1}{2}\rho v_{1}^{2} + p_{1} = \frac{1}{2}\rho v_{2}^{2} + p_{2} \implies v_{2}^{2} - v_{1}^{2} = \frac{2(p_{1} - p_{2})}{\rho}
\]
This equation has two unknowns ($v_1$ and $v_2$). We use volume conservation to eliminate $v_2$:
\[
v_{1} A_{1} = v_{2} A_{2} \implies v_{1}\left(\frac{\pi D^2}{4}\right) = v_{2}\left(\frac{\pi d^2}{4}\right) \implies v_{2} = v_{1}\left(\frac{D}{d}\right)^{2}
\]
Substitute $v_2$ back into the Bernoulli equation:
\[
v_{1}^{2}\left(\frac{D}{d}\right)^{4} - v_{1}^{2} = \frac{2(p_{1} - p_{2})}{\rho} \implies v_{1}^{2}\left[ \left(\frac{D}{d}\right)^{4} - 1 \right] = \frac{2(p_{1} - p_{2})}{\rho}
\]

\begin{resultboxnamed}{Venturi Inlet Velocity}
\[
v_{1} = \sqrt{\frac{2(p_{1} - p_{2})}{\rho \left[ \left(\frac{D}{d}\right)^{4} - 1 \right]}}
\]
\end{resultboxnamed}

\subsection{Example: Free Jet}

Consider a large tank filled with fluid up to height $H$, with a small hole of diameter $d$ at the bottom.

\begin{center}
\begin{tikzpicture}[scale=1]
    % Tank
    \draw[thick] (0, 4) -- (0, 0) -- (4, 0) -- (4, 1);
    \draw[thick] (4, 1.5) -- (4, 4);
    
    % Fluid
    \fill[primary!20] (0, 0) -- (4, 0) -- (4, 1) -- (4.2, 1) .. controls (5, 1) and (6, 0.5) .. (6, 0) -- (5.5, 0) .. controls (4.8, 0.5) and (4.2, 1.5) .. (4, 1.5) -- (4, 3.5) -- (0, 3.5) -- cycle;
    
    % Water level
    \draw[thick, blue] (0, 3.5) -- (4, 3.5);
    \node at (2, 3.7) {$\nabla$};
    
    % Points
    \filldraw[black] (2, 3.5) circle (1.5pt) node[below left] {$1$};
    \filldraw[black] (4.5, 1.25) circle (1.5pt) node[above] {$2$};
    
    % Streamline
    \draw[dashed, accent, thick, ->] (2, 3.5) .. controls (2, 1.25) and (4, 1.25) .. (4.5, 1.25);
    
    % Annotations
    \draw[<->] (-0.5, 0) -- (-0.5, 3.5) node[midway, left] {$H$};
    \draw[<->] (4, 1) -- (4, 1.5) node[below left] {$d_h$};
\end{tikzpicture}
\end{center}

Assume steady or quasi-steady flow with no shear forces. Apply Bernoulli along the streamline from the surface (1) to the exiting jet (2):
\[
\frac{1}{2}\rho v_{1}^{2} + \rho g z_{1} + p_{1} = \frac{1}{2}\rho v_{2}^{2} + \rho g z_{2} + p_{2}
\]
Assumptions:
\begin{itemize}
    \item The tank is large: $v_{1} \approx 0$
    \item Heights: $z_{1} = H$, $z_{2} = 0$
    \item Both points are exposed to the atmosphere: $p_{1} = p_{2} = p_{atm}$
\end{itemize}
\[
\rho g H = \frac{1}{2}\rho v_{2}^{2} \implies v_{2} = \sqrt{2gH} \quad \text{(Torricelli's Law)}
\]
The ideal volume flow rate is:
\[
\dot{V}_{ideal} = v_{2} A_{2} = \sqrt{2gH}\left(\frac{\pi}{4}d_{h}^{2}\right)
\]

\paragraph{Vena Contracta}
Because the streamlines must bend to exit the hole, the actual jet diameter ($d_j$) contracts to a size smaller than the physical hole diameter ($d_h$).
\begin{eqbox}{Contraction Coefficient ($C_c$)}
\[
C_{c} = \frac{A_{j}}{A_{h}} = \left(\frac{d_{j}}{d_{h}}\right)^{2}
\]
where $d_{j} \le d_{h}$. The true flow rate is scaled by $C_c$.
\end{eqbox}